\documentclass[11pt,a4paper]{article}
\usepackage[T1]{fontenc}
\usepackage[utf8]{inputenc}
\usepackage{lmodern}
\usepackage{amsmath,amssymb,amsthm,mathtools}
\usepackage[a4paper,margin=24mm,headheight=15pt,footskip=27pt]{geometry}
\usepackage{microtype,booktabs,array,longtable,tabularx}
\usepackage{xcolor,enumitem,fancyhdr,titlesec,needspace}
\usepackage[most]{tcolorbox}
\usepackage{xurl,hyperref,bookmark,listings}
\definecolor{Ink}{HTML}{182B38}
\definecolor{Accent}{HTML}{125B70}
\definecolor{Pale}{HTML}{EFF5F7}
\definecolor{Rule}{HTML}{CCDCE2}
\hypersetup{colorlinks=true,linkcolor=Accent,citecolor=Accent,urlcolor=Accent,
 pdftitle={Six Sheets, Three Escapes: Complete Fiber Geometry and Explicit Normalization of a Five-Dimensional Keller Map},
 pdfauthor={Research prepared with ChatGPT for the ProveIt project},
 pdfsubject={Exact fibers, normalization, discriminants, nonproperness, and Galois groups}}
\urlstyle{same}
\titleformat{\section}{\Large\bfseries\color{Accent}}{\thesection}{0.6em}{}
\titleformat{\subsection}{\large\bfseries\color{Ink}}{\thesubsection}{0.6em}{}
\pagestyle{fancy}\fancyhf{}
\fancyhead[L]{\small Six Sheets, Three Escapes}
\fancyhead[R]{\small ProveIt research article}
\fancyfoot[C]{\thepage}
\renewcommand{\headrulewidth}{0.35pt}
\setlength{\parskip}{4pt plus 1pt}
\setlength{\parindent}{1.1em}
\setlength{\emergencystretch}{2em}
\setlist{itemsep=3pt,topsep=5pt}
\allowdisplaybreaks[2]
\newtheorem{theorem}{Theorem}[section]
\newtheorem{lemma}[theorem]{Lemma}
\newtheorem{proposition}[theorem]{Proposition}
\newtheorem{corollary}[theorem]{Corollary}
\theoremstyle{definition}
\newtheorem{definition}[theorem]{Definition}
\newtheorem{remark}[theorem]{Remark}
\newcommand{\A}{\mathbb A}
\newcommand{\C}{\mathbb C}
\newcommand{\Q}{\mathbb Q}
\newcommand{\R}{\mathbb R}
\newcommand{\Z}{\mathbb Z}
\newcommand{\Gm}{\mathbb G_m}
\newcommand{\Spec}{\operatorname{Spec}}
\newcommand{\Disc}{\operatorname{Disc}}
\newcommand{\Res}{\operatorname{Res}}
\newcommand{\Gal}{\operatorname{Gal}}
\newcommand{\Ram}{\operatorname{Ram}}
\newcommand{\simp}{\operatorname{simp}}
\newcommand{\codim}{\operatorname{codim}}
\newcommand{\rk}{\operatorname{rank}}
\newcommand{\id}{\operatorname{id}}
\newcommand{\F}{\mathcal F}
\newcommand{\base}{\mathcal A}
\newcommand{\normal}{\mathcal N}
\newcommand{\ord}{\mathcal O}
\newcommand{\kk}{\kappa}
\newcommand{\wt}{\widetilde}
\newcommand{\code}[1]{\texttt{\nolinkurl{#1}}}
\newtcolorbox{claimbox}{breakable,colback=Pale,colframe=Rule,boxrule=0.5pt,
 arc=1mm,left=2mm,right=2mm,top=1.5mm,bottom=1.5mm}
\lstset{basicstyle=\ttfamily\small,breaklines=true,columns=fullflexible,
 keepspaces=true,showstringspaces=false,frame=single,rulecolor=\color{Rule},
 backgroundcolor=\color{Pale},aboveskip=8pt,belowskip=8pt}
\begin{document}
\hypersetup{pageanchor=false}
\begin{titlepage}
\color{Ink}
\vspace*{12mm}
{\small\sffamily\bfseries\color{Accent} ALGEBRAIC GEOMETRY \; / \; GALOIS THEORY \; / \; EXACT CERTIFICATES\par}
\vspace{17mm}
{\Huge\bfseries Six Sheets,\par Three Escapes\par}
\vspace{8mm}
{\LARGE Complete Fiber Geometry and Explicit Normalization\par}
\vspace{3mm}
{\Large of a Five-Dimensional Keller Map\par}
\vspace{14mm}
{\large A research extension of the Jacobian-conjecture\par direction in the ProveIt repository\par}
\vspace{7mm}
{24 September 2026\par}
\vfill
\begin{claimbox}
\textbf{Principal result.} We determine every geometric fiber of Gao's map
$F_6$, including the previously untreated degeneration at its first target
coordinate zero. We construct its finite normalization with an explicit
six-element free basis, identify its omitted surface and its two-component
nonproperness locus, and prove geometric and arithmetic $S_6$ results.

\textbf{Status.} The map and its generic degree are prior work. The extensions
are proved here with supplementary exact computer-algebra checks. This is an
AI-assisted research draft prepared with ChatGPT, not a peer-reviewed paper
or a Lean/Rocq-certified development. Priority beyond the inspected sources
has not been established.
\end{claimbox}
\vspace{6mm}
{\small Prepared for the ProveIt research program. Source, PDF, reproducibility
scripts, sparse coefficient data, and a verification report accompany this article.\par}
\end{titlepage}
\hypersetup{pageanchor=true}

\section*{Abstract}
\addcontentsline{toc}{section}{Abstract}
Gao's five-dimensional polynomial map $F_6$ has constant Jacobian determinant
$-290$ and generic geometric degree six. Its published construction provides
a sextic controlling fibers away from a degenerating target hyperplane,
but leaves the complete fiber stratification to further work. We resolve
the fiber-cardinality problem for this particular map and give additional
structural results.

The decisive identity is that the sweep parameter is the derivative of the
fiber sextic on its zero set. After a rescaling, this produces a uniform
root chart valid even on the exceptional hyperplane: the source open set
$x\ne0$ is exactly a simple-root locus. The remaining source hyperplane
maps isomorphically to the target hyperplane. Consequently the only
geometric fiber cardinalities are $0,1,2,3,4,6$, all of which occur. The
omitted set is an explicitly parametrized smooth closed surface isomorphic
to $\Gm\times\A^1$.

We then normalize the monogenic sextic order by adjoining one explicit
fraction. The resulting algebra is free of rank six, with three short
multiplication rules, and its spectrum is smooth. The original affine
space is the complement of the ramification locus and three unramified
boundary components in this finite model. Its nonproperness locus is
$V(cH)$, where $H$ is an irreducible normalized discriminant satisfying
$H|_{c=0}=-108(A^2+4B)$. A one-parameter Morse slice proves geometric
monodromy $S_6$; reductions modulo $3,13,37$ certify an explicit rational
specialization with the same Galois group. We also determine the possible
real fiber sizes, extract a general normalization lemma, and give an
exact inversion procedure and a proposed formalization and research agenda.

\medskip
\noindent\textbf{Keywords.} Keller map; polynomial fiber; normalization;
nonproperness; discriminant; monodromy; sextic; radical obstruction.

\newpage
\begingroup
\small
\setlength{\parskip}{0pt}
\tableofcontents
\endgroup
\newpage

\Needspace{6\baselineskip}
\section{The research target and the main conclusions}
\label{sec:intro}

\Needspace{5\baselineskip}
\subsection{Relation to ProveIt and to the published problem}
The Jacobian-conjecture project in ProveIt contains explicit polynomial
maps, exact collision certificates, symmetry arguments, and stable
reductions of degree. Its research README studies degree and sparsity
questions for the three-dimensional example. The repository walkthrough
also points to Gao's account of the subsequent higher-dimensional
constructions \cite{ProveIt,Gao}. These give a natural opportunity to pass
from individual collision witnesses to the full geometry of an explicit
map.

We focus on the map denoted $F_6$ in \cite[\S4.5.1, Theorem 4.5]{Gao}.
The subscript labels the example; its \emph{ambient dimension is five}.
The map, its determinant, its sextic elimination equation, and the
four-point fibers over its nonzero distinguished axis are prior results.
Gao's concluding discussion and Appendix A.3 explicitly leave the
complete higher-dimensional stratifications, including that of $F_6$,
unresolved in version 1. We do not claim to have constructed a new
counterexample to the Jacobian conjecture. Our object is the full fiber
geometry and finite algebraic model of this existing example.

The source audit was scoped, not repository-wide. It covered the root
README, the Jacobian project and its research README, the relevant
walkthrough search result, and the cited primary paper. The exact
repository revision and inspected file hashes are recorded in
Appendix~\ref{app:sources}. No claim is made that every research artifact
in this large repository has been searched for equivalent results.

\Needspace{5\baselineskip}
\subsection{Conventions}
Unless indicated otherwise, $k$ is a field of characteristic zero. Equations
and chart isomorphisms are defined over $\Q$ and extend to $k$. A
\emph{geometric fiber} means the fiber after passage to an algebraic
closure of its residue field, and its cardinality counts distinct points,
not multiplicities. Statements about complex nonproperness and monodromy
use $k=\C$. Real and rational fibers are treated separately.

Write a source point as $(x,y,z_1,z_2,z_3)$ and a target point as
\[
 b=(c,A,B,D,E).
\]
The use of $D$ rather than $C$ for the fourth coordinate avoids confusion
with the field $\C$. Put
\begin{equation}\label{eq:TL}
 \kk=E+B^2-AD,\qquad T=-A+2cB-c^3\kk,\qquad L=B-cD.
\end{equation}
The coordinate change from $E$ to $\kk$ is a polynomial automorphism of
target affine space.

\Needspace{5\baselineskip}
\subsection{A precise statement of what is proved}
Define
\begin{align}
 g(w)&=2w^4-2w^3-w+1=(w-1)(2w^3-1),\label{eq:g}\\
 J(w)&=g(w)+cA,\label{eq:J}\\
 P_b(w)&=w^2J(w)-cTw-c^2L,\label{eq:P}\\
 q_b(v)&=v^2J(cv)-Tv-L.\label{eq:qcompact}
\end{align}
Thus $P_b(cv)=c^2q_b(v)$. Expanded, the uniform polynomial is
\begin{align}
 q_b(v)={}&2c^4v^6-2c^3v^5-cv^3+(1+cA)v^2\notag\\
 &+(c^3\kk-2cB+A)v+cD-B.\label{eq:q}
\end{align}

\begin{theorem}[Main conclusions]\label{thm:main}
For Gao's map $\F=F_6$ the following hold.
\begin{enumerate}[label=\textup{(\roman*)}]
\item The source open set $x\ne0$ is isomorphic over target space to
\[
 \{(b,v):q_b(v)=0,\ q_b'(v)\ne0\}.
\]
The source hyperplane $x=0$ maps isomorphically onto $c=0$.
\item Every geometric fiber has
\[
 \#\F^{-1}(b)=\#\{\text{simple roots of }q_b\}+\mathbf1_{c=0}.
\]
For $c=0$ this equals $3$ when $A^2+4B\ne0$ and $1$ otherwise.
The complete set of attained cardinalities is $\{0,1,2,3,4,6\}$.
\item The omitted set over an algebraically closed field is the smooth
closed surface $M\simeq\Gm\times\A^1$ defined by
\begin{align}
 cA&=-\frac{11}{32},\notag\\
 c^3D-c^2B&=\frac{25}{128},\label{eq:M}\\
 c^4\kk-2c^2B+cA&=\frac{5}{32}.\notag
\end{align}
In particular, $\F(\A^5)=\A^5\setminus M$ on geometric points.
\item The integral closure of the target coordinate ring in the source
function field is free of rank six with basis
\[
 1,w,w^2,w^3,w^4,\zeta.
\]
Its spectrum is smooth. Its multiplication is determined by
\begin{align}
 2w^5&=c\zeta+2w^4+w^2-(1+cA)w,\notag\\
 w\zeta&=Tw+cL,\label{eq:rules}\\
 \zeta^2&=T\zeta+LJ(w).\notag
\end{align}
\item Let $H=\Disc_w(P_b)/c^2$. This is an irreducible polynomial, and
\[
 H(0,A,B,D,E)=-108(A^2+4B).
\]
The finite model's discriminant in the displayed basis is $H/256$.
The complex nonproperness locus of $\F$ is precisely $V(cH)$.
\item The geometric monodromy of the six-sheeted cover is $S_6$.
The rational target $(1,0,0,1,0)$ gives an irreducible sextic with Galois
group $S_6$ over $\Q$, so a full inverse vector there cannot be expressed
using radicals over $\Q$.
\end{enumerate}
\end{theorem}

The theorem will be proved in components. In particular, the displayed
sextic is \emph{not} itself the finite normalization: at $c=0$ it collapses
two branches to $w=0$. Ignoring this distinction gives wrong exceptional
fiber counts and an artificial discriminant component.

\Needspace{6\baselineskip}
\section{The polynomial map, without its large expansion}
\label{sec:map}

\Needspace{5\baselineskip}
\subsection{A compact definition}
We reproduce the construction needed to identify the map unambiguously
\cite[\S4.5.1]{Gao}. Set
\[
 a_1=xy,\quad a_2=x^2z_1,\quad a_3=x^3z_2,\quad a_4=x^4z_3,
\]
and apply the polynomial stage
\begin{align}
 \gamma&=1-29a_1+999a_1^2+355a_1a_2-41553a_1^3+a_4,\notag\\
 u_1&=1+27a_1-5a_2,\label{eq:stage}\\
 u_2&=a_1-12a_2+\frac{2128}{5}a_1^2+a_3,\notag\\
 u_3&=-4a_1-10a_2.\notag
\end{align}
Introduce sweep variables
\[
 w_1=\gamma u_1,\qquad w_2=\gamma^3u_2,\qquad w_3=\gamma^4u_3.
\]
Define potentials and one scalar polynomial by
\begin{align}
 G_2&=w_2+w_1w_3+w_1^2-w_1^3,\notag\\
 G_3&=2w_1G_2+w_2,\notag\\
 G_4&=w_3-G_2^2+2w_1^4-2w_1^5,\label{eq:potentials}\\
 X_1&=-w_3+2w_2-2w_1+2w_1w_3+5w_1^2-2w_1^3
                +8w_1^4-10w_1^5.\notag
\end{align}
Put
\begin{align*}
 S_1&=X_1+\gamma,& S_2&=G_2+w_1S_1,\\
 S_3&=G_3+w_1^2S_1,& S_4&=G_4+w_2S_1.
\end{align*}
With $c=x\gamma$, the map is
\begin{equation}\label{eq:map}
 \F(x,y,z_1,z_2,z_3)=
 \left(c,\frac{S_1}{c},\frac{S_2}{c^2},\frac{S_3}{c^3},\frac{S_4}{c^4}\right).
\end{equation}
All quotients in \eqref{eq:map} cancel \emph{as polynomials}. This
cancellation is part of the construction, not a restriction to $c\ne0$.

\begin{proposition}[Prior polynomiality, with an independent certificate]
\label{prop:keller}
The expression \eqref{eq:map} defines a polynomial map over $\Q$, of
coordinate degrees $(7,38,40,42,44)$, and
$\det J_{\F}=-290$. The expanded coordinates contain respectively
$6,342,421,507,904$ nonzero monomials.
\end{proposition}
\begin{proof}
Polynomiality and the degree profile are established by the stage and
weighted-divisibility construction in \cite[Theorem 4.5]{Gao}.
For an independent finite check, the accompanying script first substitutes
$(w_1,w_2,w_3)=(\gamma u_1,\gamma^3u_2,\gamma^4u_3)$ into each $S_i$
and divides each monomial by $\gamma^i$. It then substitutes
\eqref{eq:stage}, verifies that each resulting monomial has $x$-exponent
at least $i$, and divides by $x^i$. Thus no rational evaluation or
numerical cancellation is used. Counting the resulting sparse terms gives
the stated data.

For the determinant, the monomial change to the $a_i$ has determinant
$x^{10}$. The stage has determinant $-290$; its inverse is displayed
below. The change to the $w_i$ has determinant $\gamma^8$. Direct
four-variable differentiation gives
\[
 \det\frac{\partial(S_1,S_2,S_3,S_4)}
               {\partial(\gamma,w_1,w_2,w_3)}=\gamma.
\]
Padding the sweep with $c=x\gamma$ contributes another $\gamma$.
Finally the target rescaling has determinant $c^{-10}$. On
$x\gamma\ne0$ their product is
\[
 x^{10}(-290)\gamma^8\gamma^2c^{-10}=-290.
\]
Both sides are polynomials, so the identity holds everywhere. The small
stage and sweep determinants are also independently checked in the
certificate.
\end{proof}

\Needspace{5\baselineskip}
\subsection{The stage inverse and the source hyperplane}
The stage inverse is particularly useful. Given
$(\gamma,u_1,u_2,u_3)$, recover
\begin{align}
 a_1&=\frac{2u_1-2-u_3}{58},&
 a_2&=\frac{27a_1-u_1+1}{5},\notag\\
 a_3&=u_2-a_1+12a_2-\frac{2128}{5}a_1^2,&
 a_4&=\gamma-1+29a_1-999a_1^2-355a_1a_2+41553a_1^3.
 \label{eq:stageinverse}
\end{align}
Substitution proves both compositions are the identity.

The polynomial restriction of $\F$ to $x=0$ is
\begin{align}
 c&=0,\notag\\
 A&=y,\notag\\
 B&=-\frac{4636}{5}y^2+29z_1,\label{eq:section}\\
 D&=\frac{6039}{5}y^3-1041yz_1+5z_2,\notag\\
 E&=\frac{74000249}{25}y^4-\frac{92717}{5}y^2z_1
       +5yz_2-1041z_1^2-2z_3.\notag
\end{align}
This identity can be obtained by taking the coefficients of $x^i$ before
the final divisions; the accompanying polynomial construction checks all
five coordinates independently. The diagonal coefficients $1,29,5,-2$
are nonzero, so \eqref{eq:section} is a triangular polynomial isomorphism.
It provides exactly one source point on $x=0$ over every target on $c=0$.
Its determinant $-290$ also checks the global determinant normalization.

\Needspace{6\baselineskip}
\section{The derivative identity and the uniform inverse chart}
\label{sec:chart}

\Needspace{5\baselineskip}
\subsection{Eliminating the sweep variables}
For a target $b$ set
\[
 Y_1=cA,\qquad Y_2=c^2B,\qquad Y_3=c^3D,\qquad Y_4=c^4E.
\]
Write $w=w_1$. The equations $S_i=Y_i$ and
\eqref{eq:potentials} give
\begin{align}
 w_2&=Y_3-2wY_2+w^2Y_1=:V_Y(w),\label{eq:V}\\
 w_3&=Y_4+Y_2^2-Y_1Y_3+2w^5-2w^4=:W_Y(w).\label{eq:W}
\end{align}
Indeed $G_2=Y_2-wY_1$, and eliminating $G_3$ first gives
\eqref{eq:V}; eliminating $G_4$ next gives \eqref{eq:W}.
Substitution into $G_2=w_2+ww_3+w^2-w^3$ leaves
\begin{align}
 R_Y(w)={}&2w^6-2w^5-w^3+(Y_1+1)w^2\notag\\
 &+(Y_4+Y_2^2-Y_1Y_3-2Y_2+Y_1)w+Y_3-Y_2=0.
 \label{eq:RY}
\end{align}
This is Gao's sextic. Substituting the scaled target gives exactly $P_b$.

\begin{lemma}[Derivative identity]\label{lem:derivative}
As a polynomial identity in $Y$ and $w$,
\begin{equation}\label{eq:derivative}
 Y_1-X_1\bigl(w,V_Y(w),W_Y(w)\bigr)=R_Y'(w)-2R_Y(w).
\end{equation}
Consequently, on the fiber equation, $\gamma=R_Y'(w)$.
\end{lemma}
\begin{proof}
Insert \eqref{eq:V} and \eqref{eq:W} into the displayed expression for
$X_1$ in \eqref{eq:potentials} and collect coefficients in $w$. The result
is $R_Y'-2R_Y$. This is an identity before imposing $R_Y=0$, so it
also specifies the exact ideal membership, not just agreement at selected
roots. Finally $S_1=Y_1$ says $\gamma=Y_1-X_1$.
\end{proof}

For $c\ne0$, the relation $c=x\gamma$ forces $\gamma\ne0$.
Thus exactly the \emph{simple} sextic roots survive. A multiple root is
not a multiple point of the original fiber: it produces no source point.
The distinction matters because $\F$ itself is everywhere \'etale.

\Needspace{5\baselineskip}
\subsection{A rescaling that retains the exceptional fibers}
On $x\ne0$ use the regular source function
\begin{equation}\label{eq:vforward}
 v=\frac{u_1}{x}=\frac1x+27y-5xz_1.
\end{equation}
When $c\ne0$ this is $w/c$. The identities $P_b(cv)=c^2q_b(v)$ and
$P_b'(cv)=cq_b'(v)$ suggest that the derivative of $q_b$, rather than
a quotient by $c$, is the correct denominator for inversion.

\begin{theorem}[Uniform simple-root chart]\label{thm:chart}
Let
\[
 \mathcal U=\Spec k[c,A,B,D,E,v,1/q_b'(v)]/(q_b(v)).
\]
There is an isomorphism over target space
\[
 \A^5_{x\ne0}\ \simeq\ \mathcal U.
\]
In one direction use \eqref{eq:vforward}. Conversely, for a simple root
$v$ set $\delta=q_b'(v)$ and define
\begin{align}
 x&=\delta^{-1},& \gamma&=c\delta,\notag\\
 u_1&=v/\delta,&
 u_2&=(D-2vB+v^2A)/\delta^3,\label{eq:inversechart}\\
 u_3&=(\kk+2cv^5-2v^4)/\delta^4.&&\notag
\end{align}
Recover $a_i$ using \eqref{eq:stageinverse}, and then put
\begin{equation}\label{eq:sourceinverse}
 y=a_1\delta,\quad z_1=a_2\delta^2,\quad
 z_2=a_3\delta^3,\quad z_3=a_4\delta^4.
\end{equation}
No division by $c$ occurs.
\end{theorem}
\begin{proof}
First work on $c\ne0$. Lemma~\ref{lem:derivative} gives
$\gamma=P_b'(cv)=c\delta$, hence $x=1/\delta$.
Equations \eqref{eq:V}--\eqref{eq:W} give
\[
 w_2=c^3(D-2vB+v^2A),\qquad
 w_3=c^4(\kk+2cv^5-2v^4).
\]
Dividing by $\gamma^3$ and $\gamma^4$, respectively, gives
\eqref{eq:inversechart}. The stage inverse and monomial inverse now give
\eqref{eq:sourceinverse}. These steps reverse the elimination, so they
prove both compositions on $c\ne0$.

We justify extension across $c=0$, rather than assuming it. The polynomial
$q_b$ is irreducible in $k[c,A,B,D,E,v]$: as a polynomial in $E$ it is
linear with coefficient $c^3v$, and that coefficient is coprime to its
constant term. To see the coprimality, specialize $c=0$ and $v=0$:
\[
 q_b|_{c=0}=v^2+Av-B,\qquad q_b|_{v=0}=cD-B.
\]
Neither $c$ nor $v$ divides the constant term. Gauss's lemma applied to
the fraction field of the other variables proves irreducibility. Thus
$\mathcal U$ is integral and its $c\ne0$ locus is dense.

On the source open set $x\ne0$, the functions in \eqref{eq:vforward}
and $\F$ are regular. The relations $q_b(v)=0$ and
$q_b'(v)=1/x$ hold on the dense open set $c\ne0$, so they hold throughout.
Consequently the forward map lands in $\mathcal U$. All inverse formulas
are regular on $\mathcal U$, and the two composition identities extend
from its dense $c\ne0$ locus. This proves the isomorphism everywhere.
\end{proof}

\begin{corollary}[All geometric fibers]\label{cor:count}
For every target,
\begin{equation}\label{eq:count}
 \#\F^{-1}(b)=\#\{v:q_b(v)=0,\ q_b'(v)\ne0\}+\mathbf1_{c=0}.
\end{equation}
For $c=0$,
\begin{equation}\label{eq:c0count}
 \#\F^{-1}(0,A,B,D,E)=
 \begin{cases}3,&A^2+4B\ne0,\\1,&A^2+4B=0.\end{cases}
\end{equation}
\end{corollary}
\begin{proof}
Partition the source into $x\ne0$ and $x=0$. Apply
Theorem~\ref{thm:chart} to the first part and \eqref{eq:section} to the
second. At $c=0$, the polynomial is $v^2+Av-B$, with derivative $2v+A$.
A nonzero discriminant gives two simple roots. A zero discriminant gives
one double root, which is deleted; it contributes no point.
\end{proof}

\Needspace{6\baselineskip}
\section{Every fiber cardinality and the omitted surface}
\label{sec:fibers}

\Needspace{5\baselineskip}
\subsection{A finite algebraic test at every target}
For a nonzero characteristic-zero polynomial $f$ define its simple-root
factor by
\begin{equation}\label{eq:simplepart}
 G=\gcd(f,f'),\qquad S=f/G,\qquad
 \simp(f)=\frac{S}{\gcd(S,G)},
\end{equation}
with gcds and final polynomials normalized to be monic. A root of
multiplicity $m$ occurs in $G$ with multiplicity $m-1$, and in $S$ with
multiplicity one. The second gcd removes it precisely when $m\ge2$.
Hence
\[
 \#\{\text{simple geometric roots of }f\}=\deg\simp(f).
\]
Together with \eqref{eq:count}, this is an exact complete fiber-count
criterion, including singular discriminant strata. It does not require
factorization over an algebraic closure, and it is not limited to generic
points of a discriminant component.

\begin{theorem}[Cardinality spectrum]\label{thm:spectrum}
The geometric fiber cardinalities of $\F$ are exactly
\[
 \boxed{\{0,1,2,3,4,6\}.}
\]
\end{theorem}
\begin{proof}
For $c\ne0$, $P_b$ has degree six. Exactly five simple roots are
impossible: the remaining degree one would have to be another simple
root. Thus its number of simple roots belongs to
$\{0,1,2,3,4,6\}$. For $c=0$, \eqref{eq:c0count} gives only $1$ and $3$.
Table~\ref{tab:witness} supplies exact witnesses for every listed value.
For nonzero $c$ in the table, $c=1$, so $P_b=q_b$ with the variable
renamed. The displayed factorizations and elementary gcds give the
counts; the six-point row has discriminant $-993580\ne0$.
\end{proof}

\begin{table}[htbp]
\centering\small
\renewcommand{\arraystretch}{1.4}
\begin{tabularx}{\textwidth}{@{}c>{\raggedright\arraybackslash}p{0.32\textwidth}X@{}}
\toprule
Size & Target $(c,A,B,D,E)$ & Exact fiber polynomial or criterion\\
\midrule
$0$ & $(1,-11/32,0,25/128,1773/4096)$ &
$2(w^3-w^2/2-w/8-5/16)^2$\\
$1$ & $(0,0,0,0,0)$ & $q(v)=v^2$; only the section remains\\
$2$ & $(1,-11/8,0,-1/16,171/128)$ &
$(w^2+1/4)^2(2w^2-2w-1)$\\
$3$ & $(1,-8,0,-7,79)$ &
$(w-1)^3(2w^3+4w^2+6w+7)$\\
$4$ & $(1,0,0,0,0)$ & $w^2(w-1)(2w^3-1)$\\
$6$ & $(1,0,0,1,0)$ & $2w^6-2w^5-w^3+w^2+1$\\
\bottomrule
\end{tabularx}
\caption{Exact witnesses for the entire geometric cardinality spectrum.}
\label{tab:witness}
\end{table}

\Needspace{5\baselineskip}
\subsection{Classifying all sextics with no surviving root}
The empty-fiber case can be solved in closed form, not merely by the gcd
criterion. Divide $P_b$ by $2$. Its three leading nontrivial coefficient
constraints are
\begin{equation}\label{eq:fixedcoeff}
 \frac{P_b(w)}2=w^6-w^5+0w^4-\frac12w^3+	ext{terms of degree at most two}.
\end{equation}
A sextic with no simple root has multiplicity partition
$6$, $4+2$, $3+3$, or $2+2+2$.

\begin{lemma}[Unique all-multiple member]\label{lem:square}
A polynomial of the form \eqref{eq:fixedcoeff} has no simple root if and
only if
\begin{equation}\label{eq:uniquesquare}
 P_b(w)=2\left(w^3-\frac12w^2-\frac18w-\frac5{16}\right)^2.
\end{equation}
The cubic in parentheses has three distinct roots.
\end{lemma}
\begin{proof}
For multiplicity $6$, write $(w-r)^6$. The $w^5$ coefficient forces
$r=1/6$, whereas the $w^4$ coefficient is then $15r^2\ne0$.

For multiplicity $4+2$, write $(w-r)^4(w-s)^2$. The $w^5$ coefficient
forces $s=1/2-2r$. The next two required equations are
\[
 -\frac{24r^2-8r-1}{4}=0,
 \qquad \frac{8r^3+4r^2-2r+1}{2}=0.
\]
The remainder of the second polynomial on division by the first is
$(10r+23)/36$. Thus $r=-23/10$, but
$24r^2-8r-1=3609/25\ne0$, a contradiction.

For multiplicity $3+3$, write $(w^2-sw+p)^3$. Matching $w^5$ and $w^4$
gives $s=1/3$ and $p=-1/9$. Its $w^3$ coefficient is then $5/27$, not
$-1/2$.

The remaining case is a square of a monic cubic. Write
$(w^3+aw^2+bw+d)^2$ and compare $w^5,w^4,w^3$ in that order. This gives
$a=-1/2$, $b=-1/8$, $d=-5/16$, proving uniqueness. Conversely its square
has no simple root. The cubic discriminant is $-401/128\ne0$, so the
three roots are distinct and all have multiplicity exactly two in $P_b$.
\end{proof}

\begin{theorem}[Exact omitted locus]\label{thm:omitted}
Over an algebraically closed field, $\A^5\setminus\F(\A^5)$ is exactly
the surface $M$ of \eqref{eq:M}. It is closed, smooth, irreducible, and
isomorphic to $\Gm\times\A^1$. An explicit parametrization, with
$c\ne0$ and $t\in k$, is
\begin{align}
 A&=-\frac{11}{32c},& B&=\frac{t}{c^2},\notag\\
 D&=\frac{t+25/128}{c^3},&
 E&=\frac{-t^2+(53/32)t+1773/4096}{c^4}.
 \label{eq:Mparam}
\end{align}
\end{theorem}
\begin{proof}
There are no empty fibers on $c=0$ by \eqref{eq:section}. For $c\ne0$,
Lemma~\ref{lem:square} applies. The coefficients of $w^2,w,1$ in its
right-hand side are $21/32,5/32,25/128$. Comparing these with
\eqref{eq:P} gives exactly \eqref{eq:M}.

The first equation of \eqref{eq:M} forces $c$ to be invertible even
scheme-theoretically: $c(-32A/11)=1$ in the coordinate ring. With
$t=c^2B$, the other equations solve uniquely to \eqref{eq:Mparam}.
Thus the coordinate ring is $k[c,c^{-1},t]$. This proves all assertions,
including closedness in $\A^5$, dimension two, and codimension three.
\end{proof}

\begin{remark}[Geometric image versus rational image]
The equality $\F(\A^5)=\A^5\setminus M$ is a statement about geometric
points. Over $\Q$, a surviving simple root need not be rational. There are
many rational targets outside $M$ with no rational preimage; the explicit
$S_6$ specialization in Section~\ref{sec:galois} is one of them.
\end{remark}

\Needspace{6\baselineskip}
\section{An explicit finite normalization}
\label{sec:normalization}

\Needspace{5\baselineskip}
\subsection{The monogenic order is not normal}
Let
\[
 \base=k[c,A,B,D,E],\qquad
 \ord=\base[w]/(P_b(w)).
\]
Since $2$ is a unit, $\ord$ is free of rank six over $\base$.
The polynomial $P_b$ is irreducible: as a polynomial in $E$, its leading
coefficient is $c^4w$, and its constant term is divisible by neither $c$
nor $w$. Gauss's lemma gives the assertion, as in the proof of
Theorem~\ref{thm:chart}. Thus $\ord$ is a domain.

The source function field is its fraction field. Indeed, away from $c=0$
the root chart identifies $w=cv$ and gives the full source coordinates
rationally in $b,w$. Nevertheless, $\ord$ is not normal along its
collapsed locus $c=w=0$. At $c=0$ its defining equation is
\[
 P_b(w)=w^2g(w),
\]
independent of $A,B,D,E$. That single double root has forgotten the
quadratic $v^2+Av-B$ which controls the two additional source branches.

\Needspace{5\baselineskip}
\subsection{One fraction and six basis elements}
Inside $\operatorname{Frac}(\ord)$ put
\begin{equation}\label{eq:zeta}
 \zeta=\frac{wJ(w)}{c},\qquad \normal=\base[w,\zeta].
\end{equation}
The polynomial relation $P_b=0$ gives the symmetric identities
\begin{equation}\label{eq:rankone}
 wJ=c\zeta,\qquad w(\zeta-T)=cL,\qquad
 \zeta(\zeta-T)=LJ.
\end{equation}
They are exactly the multiplication rules \eqref{eq:rules}.

\begin{theorem}[Finite free normalization]\label{thm:normalization}
The algebra $\normal$ is the integral closure of both $\ord$ and $\base$
in the source function field. It is free over $\base$ with basis
\[
 \mathcal B=(1,w,w^2,w^3,w^4,\zeta).
\]
The affine variety $Z=\Spec\normal$ is smooth over $k$, and
$\pi:Z\to\A^5$ is finite flat of degree six.
\end{theorem}
\begin{proof}
\emph{Finite freeness.} The $\base$-span of $\mathcal B$ contains $1,w$
and $\zeta$. It is stable under multiplication by $w$ and by $\zeta$:
use the three rules \eqref{eq:rules}, reducing $w^5$, $w\zeta$, and
$\zeta^2$. This span is therefore the whole algebra $\normal$.
The six elements are independent, because relative to
$1,w,\ldots,w^5$ the last element has coefficient $2/c$ on $w^5$.
The degree-six irreducibility of $P_b$ gives independence over
$\operatorname{Frac}(\base)$, and hence over $\base$.
Thus $\normal$ is finite free of rank six. The same reduction and
independence show that the three displayed relations form a presentation:
any polynomial in $w,\zeta$ reduces to the basis span, and no additional
relation can remain there. In particular its elements are integral over
$\base$.

\emph{Three useful charts.} On $c\ne0$ we have
$\normal[1/c]=\ord[1/c]$. On $w\ne0$, the second relation in
\eqref{eq:rankone} gives $\zeta=T+cL/w$, so
$\normal[1/w]=\ord[1/w]$. On $J\ne0$ put $v=\zeta/J$. The first
relation gives $w=cv$, and the third gives
\[
 v^2J(cv)-Tv-L=0.
\]
Conversely $w=cv$ and $\zeta=vJ(cv)$ satisfy all three relations. Therefore
\begin{equation}\label{eq:Jchart}
 \normal[1/J]\simeq
 \base[v,1/J(cv)]/(q_b(v)).
\end{equation}
The three opens cover $Z$: a point with $c=w=0$ has $J=1$.

\emph{Smoothness on those charts.} On $c\ne0$, work in the hypersurface
$P_b=0$. Its derivative with respect to $E$ is $c^4w$. If that derivative
vanishes, then $w=0$, and the derivative with respect to $D$ is $c^3$,
which is nonzero. Thus this open is smooth. At a point with $c=0,w\ne0$,
the equation forces $g(w)=0$ and
\[
 \frac{\partial P_b}{\partial w}=w^2g'(w)\ne0,
\]
since $g=(w-1)(2w^3-1)$ has four distinct nonzero roots.
Finally, at a point with $c=w=0$, chart \eqref{eq:Jchart} applies and
$\partial q_b/\partial B=-1$. These checks prove smoothness everywhere.

A smooth integral algebra over a characteristic-zero field is normal.
The finite integral algebra $\normal$ has the same fraction field as
$\ord$. Any element of that field integral over $\ord$ is integral over
$\normal$ and therefore belongs to $\normal$. This proves that
$\normal$ is its integral closure. Since $\ord$ is itself integral over
$\base$, it is also the integral closure of $\base$ in that field.
\end{proof}

\Needspace{5\baselineskip}
\subsection{The exact fiber of the finite model at \texorpdfstring{$c=0$}{c=0}}
Reducing \eqref{eq:rankone} modulo $c$ gives
\[
 wg(w)=0,\qquad w(\zeta+A)=0,\qquad
 \zeta(\zeta+A)=Bg(w).
\]
As $g(0)=1$, the zero locus of $w$ and the zero locus of $g$ are disjoint.
The Chinese remainder decomposition is
\begin{equation}\label{eq:CRT}
 \normal/(c)\simeq
 k[A,B,D,E,\zeta]/(\zeta^2+A\zeta-B)
 \ \times\ k[A,B,D,E,w]/(g(w)).
\end{equation}
On the second factor $\zeta=-A$; on the first $w=0$.
The fiber has length two from the quadratic factor and length four from
the residual quartic factor. This accounts for all six sheets, even where
the quadratic becomes nonreduced.

\begin{theorem}[The source as an open part of its finite model]
\label{thm:boundary}
The map $\F$ factors as an open immersion followed by the finite map:
\[
 \A^5\ \xhookrightarrow{\ j\ } Z\ \xrightarrow{\ \pi\ }\A^5.
\]
Over an algebraically closed field, let $\rho_1,\rho_2,\rho_3$ be the
three roots of $2w^3=1$, and let
\[
 E_i=\{c=0,\ w=\rho_i,\ \zeta=-A\}\subset Z.
\]
Each $E_i$ is isomorphic to $\A^4$, and $\pi$ maps it isomorphically onto
$c=0$. The exact complement of the source is
\begin{equation}\label{eq:boundary}
 Z\setminus j(\A^5)=\Ram(\pi)\ \cup\ E_1\cup E_2\cup E_3.
\end{equation}
The three $E_i$ are unramified and disjoint from $\Ram(\pi)$.
\end{theorem}
\begin{proof}
On the source define $w=\gamma u_1$. The rational function $\zeta$ of
\eqref{eq:zeta} is regular there. On $x\ne0$ it equals $vJ(cv)$ by the
root chart. On $w\ne0$ it equals $T+cL/w$. These opens cover the source,
because $x=0$ implies $\gamma=u_1=w=1$. Thus they give a morphism
$j:\A^5\to Z$ over the target.

For $c\ne0$, $\pi$ is the sextic root cover. Its unramified points are
exactly the simple roots, and Theorem~\ref{thm:chart} identifies all of
these with source points. For $c=0,w=0$, decomposition
\eqref{eq:CRT} and chart \eqref{eq:Jchart} identify the unramified points
with the simple roots of $v^2+Av-B$; all of these are again source points.
The root $w=1$ at $c=0$ is unramified, and is precisely the image of the
source hyperplane under \eqref{eq:section}. The other three quartic roots
are also unramified, but cannot come from the source: $x=0$ gives $w=1$,
whereas $x\ne0$ and $c=0$ give $\gamma=0$ and $w=0$.

This proves the asserted set-theoretic image and shows that $j$ is
injective on geometric points. It lands in the \'etale locus of $\pi$.
Both its source and that locus are \'etale over the target; a morphism
between such schemes is \'etale \cite[Tag 02GH]{StacksEtale}. Geometric
injectivity makes it universally injective. An \'etale universally
injective morphism is an open immersion \cite[Tag 025G]{StacksOpen}.
Finally the four nonzero roots of $g$ are simple everywhere over $c=0$,
so all three removed components are unramified and disjoint from the
ramification locus. The description of each $E_i$ follows directly from
\eqref{eq:CRT}.
\end{proof}

The phrase ``three escapes'' refers precisely to these three missing,
unramified sheets over the generic point of $c=0$. At a generic point of
the other exceptional component, the mechanism is different: two sheets
meet at one ramification point, and that point is removed from the source.

\Needspace{6\baselineskip}
\section{The discriminant and the nonproperness locus}
\label{sec:discriminant}

\Needspace{5\baselineskip}
\subsection{Removing a spurious discriminant factor}
Let
\begin{equation}\label{eq:universaldisc}
 \Delta(a,b,d)=\Disc_w(2w^6-2w^5-w^3+aw^2+bw+d).
\end{equation}
It is unnecessary to print its full expansion. The target substitution is
\begin{equation}\label{eq:discsub}
 a=1+cA,\qquad b=-cT=c^4\kk-2c^2B+cA,
 \qquad d=-c^2L=c^3D-c^2B.
\end{equation}

\begin{theorem}[Normalized discriminant]\label{thm:H}
The expression
\begin{equation}\label{eq:H}
 H=\frac{\Delta(1+cA,-cT,-c^2L)}{c^2}
\end{equation}
is a polynomial, irreducible over every algebraically closed field of
characteristic zero, and
\begin{equation}\label{eq:H0}
 H|_{c=0}=-108(A^2+4B).
\end{equation}
In the basis $\mathcal B$ of Theorem~\ref{thm:normalization}, the trace
pairing of $\normal/\base$ has determinant $H/256$.
\end{theorem}
\begin{proof}
The discriminant of the power basis of the monic polynomial $P_b/2$ is
$\Disc(P_b)/2^{10}$. The change from that power basis to $\mathcal B$
has determinant $2/c$. A trace discriminant changes by the square of
the basis determinant. Therefore
\begin{equation}\label{eq:discbasis}
 \Disc(\mathcal B)=\frac{4}{c^2}\frac{\Disc(P_b)}{2^{10}}
                 =\frac{H}{256}.
\end{equation}
The left-hand side belongs to $\base$ because all multiplication rules
have polynomial coefficients. This proves polynomiality of $H$.

For its value at $c=0$, first assume $A^2+4B\ne0$. The two roots of $P_b$
near zero have the form $cv_++O(c^2)$ and $cv_-+O(c^2)$, where
$v_\pm$ are the two roots of $v^2+Av-B$. Their squared difference
contributes $c^2(A^2+4B)$ to first order. The other four roots tend to
the distinct roots of $g$. The leading-coefficient factors and all
cross-differences combine to $\Disc(g)=-108$: explicitly, the product
of the four roots of $g$ is $1/2$, the cross-difference contribution is
$(1/2)^4$, and the leading-coefficient ratio is $2^{10}/2^6=16$.
Thus the coefficient of $c^2$ in $\Disc(P_b)$ is
$-108(A^2+4B)$. This formal expansion is valid on a dense open set and
therefore gives the polynomial identity \eqref{eq:H0}. It also follows
by direct substitution into the exact resultant.

For irreducibility, the double-root incidence of the polynomial in
\eqref{eq:universaldisc} is parametrized by $(r,a)$:
\begin{align}
 b&=-12r^5+10r^4+3r^2-2ar,\notag\\
 d&=10r^6-8r^5-2r^3+ar^2.\label{eq:incidence}
\end{align}
Its image is exactly the discriminant hypersurface; hence that
hypersurface is irreducible. The discriminant is reduced, not a higher
power: at $(a,b,d)=(1,0,0)$ the polynomial is
$w^2(w-1)(2w^3-1)$ with one ordinary double root, and varying $d$ gives
a simple zero of the discriminant. For example the coefficient of $d$
in $\Delta(1,0,d)$ is $432\ne0$.

After inverting $c$, substitution \eqref{eq:discsub} is a coordinate
isomorphism
\[
 (c,A,B,D,\kk)\longleftrightarrow(c,B,a,b,d).
\]
Therefore $H$ is irreducible after inverting $c$. Any additional factor
before localization would have to be a power of $c$. Equation
\eqref{eq:H0} excludes such a factor. This proves irreducibility.
\end{proof}

For a finite locally free algebra in characteristic zero, the trace
pairing is nondegenerate precisely on the \'etale locus
\cite[Tag 0BJF]{StacksDisc}. Thus $V(H)$ is the branch locus of the
finite normalization, while the factor $c^2$ in the unnormalized sextic
discriminant records its nonnormal collapse. The supplied sparse file
stores $H$ with $179$ monomials in $(c,A,B,D,\kk)$.

\Needspace{5\baselineskip}
\subsection{Two components of nonproperness}
\begin{definition}
For a polynomial map $f:\C^n\to\C^n$, its nonproperness set consists
of targets $b$ for which no neighborhood $V$ of $b$ has
$f^{-1}(V)\to V$ proper. Equivalently, there exists an unbounded sequence
of source points whose images converge to $b$.
\end{definition}

\begin{theorem}[Exact nonproperness locus]\label{thm:nonproper}
For $\F:\C^5\to\C^5$,
\begin{equation}\label{eq:nonproper}
 \mathcal N_{\F}=V(cH)=V(c)\cup V(H).
\end{equation}
These are two distinct irreducible hypersurfaces. Their intersection is
\[
 c=0,\qquad A^2+4B=0.
\]
The generic fiber on $V(c)$ has three points; the generic fiber on
$V(H)$ has four points; every point of their intersection has one point
in its original fiber.
\end{theorem}
\begin{proof}
If $cH\ne0$, the sextic has six distinct roots and the original map has
six points over the target. Since $\F$ has nonzero Jacobian, choose six
disjoint local inverse neighborhoods. On a sufficiently small common
target neighborhood they supply six inverse branches. There can be no
additional branch, because every fiber has at most six points. After
shrinking to a closed ball inside this neighborhood, the preimage is
a finite union of compact continuous images of that ball. Thus $\F$ is
proper near the target.

At a target on $V(cH)$, the fiber has fewer than six points by
Corollary~\ref{cor:count}. Suppose the map were proper over a small
connected neighborhood. A proper local homeomorphism with finite fibers
is a covering onto a union of components; its fiber cardinality is
locally constant. Nearby points with $cH\ne0$ have six preimages, a
contradiction. This proves \eqref{eq:nonproper}.

Irreducibility and the intersection statement follow from
Theorem~\ref{thm:H}. At a generic point of the discriminant there is one
double root and four simple roots, so the original fiber has four
points. The hyperplane and intersection counts are
\eqref{eq:c0count}.
\end{proof}

There is an equivalent finite-model explanation. The image under $\pi$
of the ramification locus is $V(H)$, whereas the three deleted
unramified components all map to $V(c)$. Equation \eqref{eq:boundary}
therefore describes exactly which points of the finite cover have to be
added to restore properness. The omitted surface $M$ is much smaller
than the nonproperness set: points of $M$ lose every sheet, while a
generic nonproper target retains several sheets.

\Needspace{6\baselineskip}
\section{A general quadratic-contraction normalization lemma}
\label{sec:general}
The short normalization proof is not peculiar to the large coefficients
in the stage. Its algebraic core applies to a wider class of degenerating
root covers. The following lemma records that mechanism without claiming
that every such cover is a Keller map.

\begin{theorem}[Quadratic-contraction algebra]\label{thm:general}
Let $A_0$ be a domain, let $c\in A_0$ be nonzero, and let
$J(w)\in A_0[w]$ have degree $m\ge1$ and leading coefficient
$a_m\in A_0^\times$. Choose $T,L\in A_0$ and suppose
\[
 P(w)=w^2J(w)-cTw-c^2L
\]
is irreducible over $\operatorname{Frac}(A_0)$. In its root field define
$\zeta=wJ(w)/c$ and $N=A_0[w,\zeta]$. Then:
\begin{enumerate}[label=\textup{(\roman*)}]
\item $N$ is finite free of rank $m+2$ over $A_0$, with basis
$1,w,\ldots,w^m,\zeta$ and relations
\[
 wJ=c\zeta,\qquad w\zeta=Tw+cL,\qquad
 \zeta^2=T\zeta+LJ.
\]
\item Its trace discriminant in that basis equals
\[
 \frac{\Disc(P)}{a_m^{2m}c^2}.
\]
\item On $c\ne0$ or $w\ne0$ the algebra equals the original monogenic
order. On $J\ne0$ it has the chart
\[
 A_0[v,1/J(cv)]/(v^2J(cv)-Tv-L).
\]
If $c$ and $J(0)$ generate the unit ideal, these three opens cover
$\Spec N$. If their rings are normal, $N$ is the integral closure of
the monogenic order in its fraction field.
\end{enumerate}
\end{theorem}
\begin{proof}
The three relations follow by multiplying $P=0$ by suitable powers of
$J/c$ and cancelling in the fraction field. They show that the span of
the displayed basis is closed under multiplication by $w$ and $\zeta$;
the leading coefficient $a_m$ is a unit, so $w^{m+1}$ reduces using
$wJ=c\zeta$. Independence follows from the coefficient $a_m/c$ of
$w^{m+1}$ in $\zeta$. This proves (i).

The monic power-basis discriminant is
$\Disc(P)/a_m^{2m+2}$. The basis determinant is $a_m/c$, giving (ii).
The chart calculations are the same elementary substitutions as in
\eqref{eq:Jchart}. If $c=w=J=0$, then $c=J(0)=0$, which is excluded by
the unit-ideal assumption. Thus the charts cover. Normality is local;
a finite normal algebra in the same fraction field is the integral
closure. This proves (iii).
\end{proof}

For $F_6$, take $m=4$, $a_m=2$, and $J(0)=1+cA$.
The smoothness checks in Theorem~\ref{thm:normalization} supply the
normality hypotheses. For a general $J,T,L$, those hypotheses must be
checked: the finite free algebra in this lemma is not automatically
normal. This distinction is important when attempting degree-twelve or
higher analogues.

The lemma explains the otherwise mysterious square factor $c^2$ in a
root discriminant: one basis vector must be divided by $c$ to separate
the two collapsed directions. It also provides a practical alternative
to normalizing a large graph ideal in all source and target variables.

\Needspace{6\baselineskip}
\section{Geometric monodromy and an arithmetic radical obstruction}
\label{sec:galois}

\Needspace{5\baselineskip}
\subsection{A Morse slice with full symmetric monodromy}
Set
\[
 p(w)=2w^6-2w^5-w^3+w^2.
\]
On the target line $(c,A,B,D,E)=(1,0,0,t,0)$, the fiber polynomial is
$p(w)+t$. Its discriminant is
\begin{equation}\label{eq:morsedisc}
 \Disc_w(p(w)+t)
 =-4t(32t+3)(11664t^3-5680t^2+1149t-36).
\end{equation}
The polynomial on the right is squarefree. This can be checked by the
Euclidean algorithm, and is independently verified in the code.

\begin{theorem}[Geometric $S_6$]\label{thm:geometricS6}
The six-sheeted cover over $\C^5\setminus V(cH)$ has monodromy group
$S_6$. Equivalently, $P_b$ has Galois group $S_6$ over
$\C(c,A,B,D,E)$.
\end{theorem}
\begin{proof}
For a polynomial $p$ of degree six, the discriminant of $p(w)+t$ is a
nonzero constant times the product of $t+p(r)$ over the five critical
points $r$ of $p$, counted with multiplicity. The squarefreeness of
\eqref{eq:morsedisc} implies that all five critical points are simple
and their critical values are distinct.

Remove the five critical values from the target $t$-line. The inverse
image under $t=-p(w)$ is the complex $w$-line with finitely many points
removed, hence is connected. Its degree-six covering therefore has
transitive monodromy. A small loop around a simple critical value
interchanges exactly two sheets, because the local map has a nonzero
quadratic term. The fundamental group of the punctured line is generated
by these small loops. Thus its monodromy is a transitive permutation
group generated by transpositions.

Place an edge between two sheet labels whenever a generating
transposition exchanges them. Transitivity makes this graph connected.
Transpositions along a connected graph generate all permutations, so
the group is $S_6$. Restricting the full target cover to this line cannot
enlarge its monodromy. Since the restricted group is already $S_6$, the
full group is exactly $S_6$.
\end{proof}

\Needspace{5\baselineskip}
\subsection{An explicit rational certificate}
Consider the rational target
\[
 b_*=(1,0,0,1,0),\qquad
 f(w)=P_{b_*}(w)=2w^6-2w^5-w^3+w^2+1.
\]
Its discriminant is $-993580$. The following factorizations are in the
indicated finite fields:
\begin{align}
 f&\equiv2(w^6-w^5+w^3-w^2-1)&&\pmod3,\label{eq:mod3}\\
 f&\equiv2(w-4)(w^5+3w^4-w^3+2w^2+2w-5)&&\pmod{13},\label{eq:mod13}\\
 f&\equiv2(w+14)(w-16)(w-7)(w-5)(w^2+13w+14)&&\pmod{37}.
 \label{eq:mod37}
\end{align}
The sextic in \eqref{eq:mod3}, the quintic in \eqref{eq:mod13}, and the
quadratic in \eqref{eq:mod37} are irreducible in their respective fields.
The supplied verifier proves this with modular exponentiation and gcds,
not by accepting a black-box Galois-group output. Specifically, for a
monic polynomial $h$ of degree $n$ over $\mathbb F_p$, it checks
\[
 X^{p^n}\equiv X\pmod h,
 \qquad
 \gcd(h,X^{p^{n/\ell}}-X)=1
 \quad(\ell\text{ prime},\ \ell\mid n).
\]
These conditions are equivalent to irreducibility, by the factorization
of $X^{p^r}-X$ into the monic irreducibles whose degrees divide $r$.

\begin{theorem}[Arithmetic $S_6$ and inverse obstruction]
\label{thm:arithmeticS6}
The polynomial $f$ is irreducible over $\Q$ and has Galois group $S_6$.
No one of its roots is expressible by radicals over $\Q$. In particular,
no full source point of $\F^{-1}(b_*)$ has all five coordinates expressible
by radicals over $\Q$.
\end{theorem}
\begin{proof}
Irreducibility modulo $3$ implies irreducibility over $\Q$, and hence a
transitive Galois action. At the three displayed primes the reductions
are squarefree and the leading coefficient remains a unit. The
factorization-cycle theorem gives a $5$-cycle from $13$ and a
transposition from $37$ \cite[Theorem 4.28]{Milne}. To fit the monic
integer formulation exactly, apply it to $2^5f(W/2)$; at these odd
primes the change of variable preserves factor degrees.

A transitive subgroup of $S_6$ containing a $5$-cycle has a point
stabilizer transitive on the other five labels, and is therefore
$2$-transitive. Conjugating the transposition gives every transposition.
Hence the group is $S_6$, in agreement with the standard group criterion
\cite[Lemma 4.31]{Milne}.

An irreducible polynomial with nonsolvable splitting group cannot have
a root expressible by radicals: all conjugates of such a root would
belong to a common solvable normal radical extension after adjoining
suitable roots of unity, contradicting nonsolvability of $S_6$
\cite[\S5, solvability theorem]{Milne}.
Finally $w=\gamma u_1$ is a polynomial in the source coordinates with
rational coefficients. If an entire inverse vector had radical
coordinates, its corresponding root $w$ would be radical as well.
\end{proof}

This result does not forbid an exact inverse expressed by algebraic
numbers or by a root-isolation representation. It forbids a radical-only
inverse vector in this explicit fiber. It is therefore directly relevant
to the repository's distinction between exact algebraic solving and
radical-solvability certificates.

\Needspace{6\baselineskip}
\section{Real fibers and arithmetic inversion}
\label{sec:real}

\begin{theorem}[Real cardinalities]\label{thm:real}
For the real polynomial map $\F:\R^5\to\R^5$, the possible fiber
cardinalities are exactly $\{0,1,2,3,4\}$. If $c\ne0$, the possible
values are contained in $\{0,1,2,4\}$. If $c=0$, the count is $3$ for
$A^2+4B>0$ and $1$ for $A^2+4B\le0$.
\end{theorem}
\begin{proof}
For $c\ne0$ the real source points correspond to the real simple roots
of $P_b$. Suppose all six roots $r_i$ were real, counted with
multiplicity. Newton's identities applied to \eqref{eq:fixedcoeff} give
\[
 \sum_i r_i=1,\qquad \sum_i r_i^2=1,\qquad
 \sum_i r_i^3=\frac52.
\]
But
\[
 \left|\sum_i r_i^3\right|\le\sum_i|r_i|^3
 \le\left(\sum_i r_i^2\right)^{3/2}=1,
\]
a contradiction. Five simple real roots would force the remaining root
to be real as well. Exactly three simple real roots would leave a
real cubic factor with no additional simple real root. Its odd degree
forces a real root, and that root must be triple; otherwise a further
simple real root remains. This would again make all six roots real.
Thus only $0,1,2,4$ simple real roots can occur.

For $c=0$, use $q_b=v^2+Av-B$ and the one-point section. Two distinct
real roots give three points; a double root is deleted, and nonreal
roots give no real points. This proves the stated criterion.

All five global values occur. Zero occurs at $b_*$: write
$p(w)=w^2(w-1)(2w^3-1)$. It is nonnegative outside
$(2^{-1/3},1)$, while on that interval $-1<p(w)<0$. Thus $p(w)+1>0$
on all of $\R$. One occurs at the origin, two at $(1,0,0,0,0)$, and
three at $(0,0,1,0,0)$. Four occurs at $(1,-8,0,-6,70)$, whose sextic is
\[
 2w^6-2w^5-w^3-7w^2+14w-6.
\]
It has the simple root $1$, and sign changes on
$(-2,-1)$, $(0,9/10)$, and $(11/10,2)$. These give four distinct real
roots; the preceding bound proves there are no others. Exact Sturm
counts independently verify all five examples.
\end{proof}

A particularly instructive target is therefore $b_*$. Its geometric
fiber has six points, its real fiber is empty, and its rational fiber is
empty. These statements coexist without contradiction: geometric
surjectivity onto $\A^5\setminus M$ does not assert real or rational
surjectivity.

\Needspace{5\baselineskip}
\subsection{An exact inversion procedure over a computable field}
The results give a compact algorithm which never expands or eliminates
all $2180$ monomials of the original map.
\begin{enumerate}[label=\textbf{Step \arabic*.}]
\item Form $\kk,T,L$ and $q_b$ by \eqref{eq:TL}--\eqref{eq:q}. For geometric
cardinality, compute the gcds in \eqref{eq:simplepart} and take the degree,
adding one when $c=0$.
\item For inverse points, factor or isolate the roots of $\simp(q_b)$ in
the desired field. Over $\Q$, degree-one factors are exactly the rational
roots; higher-degree factors represent algebraic inverse points.
\item At each simple root evaluate \eqref{eq:inversechart},
\eqref{eq:stageinverse}, and \eqref{eq:sourceinverse}. The only variable
denominator is a power of $q_b'(v)$, whose nonvanishing has already been
certified.
\item If $c=0$, also solve the triangular equations \eqref{eq:section}.
This supplies the unique point not represented by the $x\ne0$ chart.
\end{enumerate}
Polynomial degree is at most six throughout the root computation, and
drops to two on $c=0$. The formula handles multiplicities exactly; a
near-zero numerical derivative must never be substituted for a gcd test.
The accompanying Python module exposes the root polynomial, simple-part
calculation, rational reconstruction, and section inverse as reusable
functions.

\Needspace{6\baselineskip}
\section{Verification, proof boundaries, and formalization}
\label{sec:verification}

\Needspace{5\baselineskip}
\subsection{What the reproducibility package checks}
The package contains a sparse rational construction of $\F$, a separate
verification driver, two coefficient files, and a saved verification log.
The driver completed \textbf{63 exact checks} under Python 3.13.5 and
SymPy 1.14.0. No floating-point arithmetic or network access is used.

The checks include cancellation of the apparent denominators, degree and
term counts, the entire $x=0$ section, the stage inverse, the stage and
sweep determinants, the derivative identity, the rescaled root equation,
the three normalization relations, the exceptional quadratic, the exact
$c^2$ discriminant factor and its leading coefficient, exclusion of the
other all-multiple partitions, every cardinality witness, exact inverse
round trips including $\gamma=0$, the Morse discriminant, finite-field
irreducibility certificates, and real Sturm counts.

The saved log is a reproducibility record, not an independent proof
assistant certificate. In particular, a finite selection of inverse
round trips cannot establish a chart isomorphism. That is why
Theorem~\ref{thm:chart} proves the identities on a dense open set and
then explicitly proves that they extend over the exceptional locus.
Likewise, normality, openness of the embedding, and the monodromy
argument are proved mathematically; they are not inferred from a
computer-algebra ``PASS'' message.

\Needspace{5\baselineskip}
\subsection{A suitable Lean development boundary}
A formalization should begin with the small algebra, not with a giant
five-variable graph ideal. The natural dependency order is:
\begin{enumerate}[label=\textup{(\arabic*)}]
\item Prove the elimination and derivative identities in a polynomial
ring. These are finite ring equalities with rational coefficients.
\item Formalize the multiplicity-one factor construction
\eqref{eq:simplepart}, including the zero-characteristic hypotheses and
constant-polynomial edge case. This yields a reusable exact fiber-count
checker.
\item Establish the six-element algebra and its multiplication rules.
Prove spanning by closure under the two generators and independence by
the degree-six field relation. Then derive the discriminant basis change.
\item Formalize the source chart and its exceptional section, either as
localized coordinate-ring isomorphisms or first as field-valued point
bijections with explicit nonvanishing hypotheses.
\item Add the geometric layer: normality of the three charts, the open
immersion criterion, and the branch/nonproperness interpretation. This
stage needs substantially more algebraic-geometric infrastructure than
the finite polynomial certificates.
\end{enumerate}
No Lean source is supplied as if it were checked. A future kernel audit
should report separately which finite identities, pointwise statements,
and scheme-theoretic conclusions have actually been formalized.

\Needspace{5\baselineskip}
\subsection{Why the argument is not a numerical discovery claim}
The important transitions are algebraic. A root is deleted because an
exact derivative is zero. An entire fiber is empty because an exhaustive
list of multiplicity partitions reduces to one square. The finite model
is normal because its charts are smooth, not because sample fibers look
correct. The Galois group is $S_6$ because of exact permutations forced
by finite-field factorizations and a connected Morse cover. These
arguments continue to hold at targets which numerical root solvers
would find highly ill-conditioned.

\Needspace{6\baselineskip}
\section{Further research questions and proposed projects}
\label{sec:future}
The following questions are specific continuations of the results above.
They are proposed directions, not additional theorems established here.

\Needspace{5\baselineskip}
\subsection{Complete the adjacent higher-dimensional fiber problems}
\textbf{1. The degree-twelve map $F_7$.} Gao supplies a univariate
fiber polynomial of degree twelve \cite[\S4.6]{Gao}. Determine whether
its exceptional degeneration is resolved by successive versions of
Theorem~\ref{thm:general}. A concrete first deliverable is an integral
basis and a factorization of its normalized discriminant, followed by
an exhaustive empty-fiber classification. The $F_6$ analysis suggests
that the exceptional hyperplane must be treated independently of the
generic elimination equation.

\textbf{2. The maps $F_4$ and $F_5$.} Their published quintic tangency
and degree-ten resultant descriptions leave related fiber problems
\cite[Appendix A.3]{Gao}. Derive a derivative or subresultant criterion
that distinguishes an actual source point from a deleted tangency
parameter. The main risk is that a resultant introduces extraneous
solutions; each proposed root representation should be accompanied by
an explicit inverse chart.

\Needspace{5\baselineskip}
\subsection{Refine the geometry of the sextic example}
\textbf{3. Singular strata of $H$.} The simple-root gcd criterion gives
the correct count at every point, but does not constitute a Whitney
stratification. Determine the singular locus of the discriminant via
\eqref{eq:incidence}, separating triple roots, two double roots, and
higher coincidences. Describe their intersections with $c=0$ and the
omitted surface $M$. An attractive deliverable is a small set of
parametrizations or subresultant ideals rather than a large expanded
elimination basis.

\textbf{4. Quantitative escape rates.} For each deleted component $E_i$,
construct explicit Laurent-series source arcs whose images tend to a
prescribed point of $c=0$. Compute the leading pole orders of all five
source coordinates. Repeat near a generic ramification point and near
the triple-double locus $M$. This would distinguish the unramified and
ramified mechanisms dynamically, not just set-theoretically.

\textbf{5. The smooth affine normalization $Z$.} Determine its divisor
class group, Picard group, and algebraic automorphisms. Is the explicit
six-generator order monogenic after a target automorphism or after a
smaller localization? Any answer must respect the quadratic-plus-quartic
special fiber \eqref{eq:CRT}; the naive power basis cannot do so.

\Needspace{5\baselineskip}
\subsection{Extend the algebraic and arithmetic methods}
\textbf{6. Higher contractions.} Replace $w^2J-cTw-c^2L$ by a family
whose leading degeneration contains $w^r$, with $r\ge3$. Find hypotheses
under which $r-1$ explicit divided basis elements suffice to normalize
it. The desired theorem should state a finite-free basis, multiplication
rules, discriminant correction, and checkable normality conditions.

\textbf{7. Arithmetic image and solvable fibers.} Classify rational
parameter subvarieties on which the sextic Galois group drops to a
solvable subgroup. Distinguish those with a rational source point from
those with only nonrational radical inverse points. The explicit $S_6$
certificate is a starting obstruction, not a classification of all
specializations. This problem connects directly to ProveIt's
radical-solvability and exact polynomial-solving projects.

\textbf{8. Good reduction.} Analyze the root chart and normalization in
positive characteristic, separating primes excluded by the map's
coefficients or determinant from primes affecting discriminants such as
$401$. The characteristic-zero proof must not simply be reduced modulo
an arbitrary prime: derivative multiplicities, separability, and the
stage determinant require independent checks.

\Needspace{5\baselineskip}
\subsection{Formal and structural questions}
\textbf{9. A certified root-chart toolkit.} Build a reusable library that
takes a polynomial identity for a graph, a reconstruction tuple, and
nonvanishing certificates, and returns a checked fiber bijection. Add
support for exceptional source components instead of silently localizing
them away. This would serve the Jacobian, radical-expression, and
algebraic-number parts of ProveIt simultaneously.

\textbf{10. Rigidity versus finite-cover complexity.} Gao's middle-branch
rigidity questions concern when a higher-dimensional construction is
only a suspension of a lower-dimensional one \cite[\S4.5.2]{Gao}.
Investigate whether normalized-cover invariants, especially monodromy
and the number of unramified boundary components, obstruct such
suspensions. The present $S_6$ computation alone is not a proof of a
new rigidity theorem; a comparison theorem for finite normalizations
would be needed.

\Needspace{6\baselineskip}
\section{Conclusion}
The central open task selected from the repository's surrounding research
program was not the existence of another collision, but the exact
behavior of every fiber of an already explicit high-degree map.
The derivative identity makes that task accessible. The uniform chart
then exposes the exceptional quadratic, and one divided basis element
repairs the finite model. These steps determine the image, all geometric
and real cardinalities, the branch and nonproperness loci, and a concrete
radical obstruction.

The resulting model is small: a sextic, a quartic, a quadratic exceptional
fiber, and three multiplication rules. It retains the geometry hidden by
the large expanded map and provides a practical target for independent
review, further generalization, and formal verification.

\appendix
\Needspace{6\baselineskip}
\section{Reproducing the calculations}
\label{app:reproduce}
From the extracted archive directory, run
\begin{lstlisting}
python -m pip install -r requirements.txt
python code/verify.py
pdflatex -interaction=nonstopmode -halt-on-error Six_Sheets_Three_Escapes.tex
pdflatex -interaction=nonstopmode -halt-on-error Six_Sheets_Three_Escapes.tex
\end{lstlisting}
The Python scripts do not require a LaTeX installation. The article
compiles independently of Python: the mathematical expressions and
proofs are contained directly in the TeX source.

\begin{center}\small
\renewcommand{\arraystretch}{1.25}
\begin{tabularx}{\textwidth}{@{}>{\raggedright\arraybackslash}p{0.43\textwidth}X@{}}
\toprule
File & Purpose\\
\midrule
\code{code/build_map.py} & Sparse rational construction and exact denominator cancellation.\\
\code{code/verify.py} & The 63-check certificate and reusable rational inverse functions.\\
\code{data/map_coefficients.json} & All five polynomial coordinates; monomial exponent order is $(x,y,z_1,z_2,z_3)$.\\
\code{data/normalized_discriminant.json} & Sparse $H$ in coordinate order $(c,A,B,D,\kk)$; $\kk=E+B^2-AD$.\\
\code{data/verification_report.txt} & Saved output of the completed exact run.\\
\code{SOURCE_AUDIT.md} & Repository revision, source scope, and novelty boundary.\\
\bottomrule
\end{tabularx}
\end{center}

Every coefficient in the JSON files is a rational string; no precision
parameter is needed. The verifier regenerates the coefficient files
from the formulas rather than accepting them as trusted input.
A failure raises an exception and prevents the final all-checks-passed
message. There are no random checks in the verifier.

\Needspace{6\baselineskip}
\section{Source audit and attribution}
\label{app:sources}
The repository reference observed during the audit was
\begin{center}\small
\code{e21766d04c2b8a9b2cdba0cd43e563024b3bd1b9}.
\end{center}
The main directly inspected project files were
\begin{center}\small
\code{Algebra/JacobianConjecture/README.md}\\
\code{Algebra/JacobianConjecture/Research/README.md}.
\end{center}
Their returned blob hashes were respectively
\begin{center}\small
\code{67e25dcb94cb7a38fb508ea65e3a35455df859ac}\\
\code{f0605a788b98bed23df91318aa2069e6802fc66e}.
\end{center}
The root README and a walkthrough search result were used for orientation.
The root README's general formalization claims were not independently
rebuilt in this work, and none is used as a substitute for the specific
algebraic proofs above.

The main external mathematical source was Gao's version 1, dated
31 July 2026, particularly \S4.5.1, Theorem 4.5, the concluding remarks,
and Appendix A.3. Its definition of $F_6$, determinant, generic degree,
and elimination polynomial are explicitly credited. The derivative
identity, uniform chart, complete cardinality criterion, exact omitted
surface, six-element normalization, normalized discriminant and boundary,
and the geometric and arithmetic $S_6$ proofs are the extensions developed
in this article. This describes their relation to the inspected sources,
not a guarantee that no equivalent result exists elsewhere.

The primary literature search checked the cited paper's available version
and searches combining its identifier or author with $F_6$, fiber,
normalization, and Jacobian terminology. No directly matching extension
was identified in that bounded search. It was not an exhaustive review
of all unpublished work, repository branches, or later inaccessible
material. The Stacks Project and Milne are cited for the standard
geometric and Galois-theoretic facts used in the proofs.

\begin{thebibliography}{9}
\addcontentsline{toc}{section}{References}
\bibitem{ProveIt}
Vladimir Reshetnikov, \emph{ProveIt}, GitHub repository, Jacobian-conjecture
project and research README; revision
\code{e21766d04c2b8a9b2cdba0cd43e563024b3bd1b9}.
\url{https://github.com/VladimirReshetnikov/ProveIt}.

\bibitem{Gao}
Shuhong Gao, \emph{Counterexamples to the Jacobian conjecture in dimensions
greater than two}, arXiv:2608.00222v1, 31 July 2026.
\url{https://arxiv.org/abs/2608.00222v1}.
In particular \S4.5.1, Theorem 4.5, and Appendix A.3.

\bibitem{StacksEtale}
The Stacks Project Authors, \emph{The Stacks Project}, Tag 02GH,
``\'Etale morphisms,'' including the lemma that a morphism between
schemes \'etale over a common base is \'etale.
\url{https://stacks.math.columbia.edu/tag/02GH}.

\bibitem{StacksOpen}
The Stacks Project Authors, \emph{The Stacks Project}, Tag 025G,
Theorem 41.14.1: universally injective \'etale morphisms are open immersions.
\url{https://stacks.math.columbia.edu/tag/025G}.

\bibitem{StacksDisc}
The Stacks Project Authors, \emph{The Stacks Project}, Tag 0BJF,
Lemma 49.3.1: the discriminant of a finite locally free morphism and
its \'etale locus.
\url{https://stacks.math.columbia.edu/tag/0BJF}.

\bibitem{Milne}
James S. Milne, \emph{Fields and Galois Theory}, version 5.10,
September 2022, 144 pp. Theorem 4.28, Lemma 4.31, and Chapter 5.
\url{https://www.jmilne.org/math/}.
\end{thebibliography}
\end{document}
